\section{Results supported after label correction}
\label{sec:findings}
\subsection{The comparison population changes the apparent detector quality}
% claims: C05 C06 C07 C08
\begin{center}\input{\evidenceroot/tables/recovery_detector.tex}\end{center}
The main-adapter day ROC changes from \rv{r42.all_users.day_roc} to
\rv{r42.unseen_users.day_roc} on r4.2 and from \rv{r62.all_users.day_roc} to
\rv{r62.unseen_users.day_roc} on r6.2 when negatives are restricted to unseen
users. This is a population/exposure comparison, not a controlled intervention
that identifies the sole cause of the gap. In r4.2, nine of the ten insiders
newly included by the label repair are scenario-3 users, and their full scores
rank incident-window days poorly. Their inclusion materially changes the main
result; recovery cannot stop at the previously covered 60 users.

The historical Deep SVDD, recurrent-autoencoder, and Isolation Forest rows are
not included as corrected comparisons. Their fitting populations need repair.
The corrected LM measurements do not establish a detector leaderboard ordering.

\subsection{Behavior-target scoring improves the complete r4.2 evaluation}
% claims: C32
\begin{center}\input{\evidenceroot/tables/recovery_behavior.tex}\end{center}
This separate, batch-size-one score decomposition covers all 42,468 unseen-pool
days, all 1,362 positive days, and all 70 insiders. A 256-example bridge exactly
reproduced the existing class-score pass before the missing users were scored.
Class losses reconstruct the total score, and an independent check reproduced
840 reported fold metrics to floating-point precision.

With the model and input held fixed, changing the average from all targets to
behavior targets raises mean fold day ROC from \rv{c32.full.day_roc} to
\rv{c32.behavior_only.day_roc}. The paired gain is \rv{c32.gain}, with a
descriptive 95\% user-bootstrap interval \rv{c32.gain.ci}. AP changes from
\rv{c32.full.day_ap} to \rv{c32.behavior_only.day_ap}. This supports a
\emph{score mitigation} for this model and population. It does not show that
profile information has stopped influencing behavior predictions.

Mean fold prevalence is \rv{c32.fold.prevalence}; pooled prevalence is
\rv{c32.pooled.prevalence}. The fold AP belongs with the former evaluation
construction, not the pooled prevalence. Neither AP value is the percentage
of alerts that would be correct at an operational threshold.
\begin{figure}[ht]
\centering\includegraphics[width=.94\textwidth]{\evidenceroot/figures/recovery_score_mitigation.pdf}
\caption{Complete r4.2 score decomposition. All three views retain the same input
and frozen model. The figure reports fold means, not pooled metrics.}
\label{fig:mitigation}
\end{figure}

\subsection{Training changes and score changes must stay separate}
% claims: C03 C09 C28 C31
The next table concerns separate r6.2 factorial adapters, not the main adapters
in the first table. Each input condition changes both adaptation and evaluation
content, with one adapter seed per condition.
\begin{center}\input{\evidenceroot/tables/recovery_factorial.tex}\end{center}
Removing all profile content improves ranking in both audited model families.
The 3B change is substantial in point estimate; the earlier description of a
uniformly tiny 3B effect is not supported. These are run-specific comparisons
on five insiders, not a general model-capacity law. Their training-byte
provenance is weaker than the main-adapter audit and remains explicitly qualified.

Scoring behavior targets while retaining the full input yields another comparison:
\begin{center}\input{\evidenceroot/tables/recovery_r62_behavior.tex}\end{center}
The 3B target-mask adapter omits profile targets from its training objective
while retaining profile context. Relative to behavior scoring of its full-training
counterpart, its day-ROC gain is about 0.006, with a descriptive interval crossing
zero, while AP decreases. This does not establish an additional mitigation gain.

Input masking of frozen main adapters is yet another experiment. Removing
personality input improves the cached day-ROC point estimates, but removing
all profile input does not have the same effect. Those arms used different
batch/hardware/cross-entropy configurations; the large personality-removal
contrast exceeds the observed metric sensitivity, while small differences
cannot be read at fine numerical precision. These findings are compatible with
several mechanisms and do not identify one by themselves.

\subsection{Temporal models provide a bounded comparison}
% claims: C33 C34
The corrected cached numeric-model comparison uses a common population of
141,662 history-eligible unseen-user days:
\begin{center}\input{\evidenceroot/tables/recovery_timeseries.tex}\end{center}
These are \emph{pooled} metrics. The seven-observation reconstruction includes
the target day; the forecast uses preceding observations only. Histories are
previous available observations, not necessarily consecutive calendar days.
All three checkpoint reloads passed the declared numerical tolerances, although
the seven-observation profile score used 99.5\% of its tolerance. That is a
tolerance pass, not bitwise equality. History-shuffle and broader copying or
mechanistic claims remain pending.
